\subsection*{2.4 The Gradient Vector}

\begin{conceptbox}
    The gradient is a vector of partial derivatives:
    \begin{equation*}
        \nabla f = \left\langle \frac{\partial f}{\partial x}, \frac{\partial f}{\partial y}, \frac{\partial f}{\partial z} \right\rangle
    \end{equation*}
\end{conceptbox}

\begin{examplebox}
    Let the function be:
    \begin{equation*}
        f(x,y,z) = x^2 + y^2 + z^2
    \end{equation*}

    ---

    \textbf{Step 1: Recall the definition of the gradient}

    The gradient of a scalar function is:
    \begin{equation*}
        \nabla f = \left\langle \frac{\partial f}{\partial x}, \frac{\partial f}{\partial y}, \frac{\partial f}{\partial z} \right\rangle
    \end{equation*}

    So we compute each partial derivative one at a time.

    ---

    \textbf{Step 2: Compute partial derivatives}

    \textbf{Partial with respect to $x$:}
    \begin{align*}
        \frac{\partial f}{\partial x}
         & = \frac{\partial}{\partial x}(x^2 + y^2 + z^2) \\
         & = 2x + 0 + 0                                   \\
         & = 2x
    \end{align*}

    \textbf{Partial with respect to $y$:}
    \begin{align*}
        \frac{\partial f}{\partial y}
         & = \frac{\partial}{\partial y}(x^2 + y^2 + z^2) \\
         & = 0 + 2y + 0                                   \\
         & = 2y
    \end{align*}

    \textbf{Partial with respect to $z$:}
    \begin{align*}
        \frac{\partial f}{\partial z}
         & = \frac{\partial}{\partial z}(x^2 + y^2 + z^2) \\
         & = 0 + 0 + 2z                                   \\
         & = 2z
    \end{align*}

    ---

    \textbf{Step 3: Form the gradient vector}

    \begin{align*}
        \nabla f = \langle 2x, 2y, 2z \rangle
    \end{align*}

    ---

    \textbf{Step 4: Evaluate at the point $(1,2,3)$}

    Substitute $x=1$, $y=2$, $z=3$:

    \begin{align*}
        \nabla f(1,2,3)
         & = \langle 2(1), 2(2), 2(3) \rangle \\
         & = \langle 2, 4, 6 \rangle
    \end{align*}

    ---

    \textbf{Step 5: Compute the magnitude of the gradient}

    The magnitude of a vector $\langle a,b,c \rangle$ is:
    \begin{equation*}
        |\mathbf{v}| = \sqrt{a^2 + b^2 + c^2}
    \end{equation*}

    Apply this to $\nabla f(1,2,3) = \langle 2,4,6 \rangle$:

    \begin{align*}
        |\nabla f|
         & = \sqrt{2^2 + 4^2 + 6^2} \\
         & = \sqrt{4 + 16 + 36}     \\
         & = \sqrt{56}
    \end{align*}

    ---

    \textbf{Final Answers:}
    \begin{align*}
        \nabla f          & = \langle 2x,2y,2z \rangle \\
        \nabla f(1,2,3)   & = \langle 2,4,6 \rangle    \\
        |\nabla f(1,2,3)| & = \sqrt{56}
    \end{align*}

    ---

    \textbf{Geometric Interpretation:}
    \begin{itemize}
        \item The gradient vector points in the direction of \textbf{maximum increase} of the function.
        \item At the point $(1,2,3)$, the function increases fastest in the direction $\langle 2,4,6 \rangle$.
        \item The magnitude $\sqrt{56}$ tells us \textbf{how fast} the function increases in that direction.
    \end{itemize}

\end{examplebox}
